\section{The result and the plan}\label{sec:intro}

\subsection{The result}

Throughout, $\log$ is the binary logarithm, and $\lgk{j}$ is its $j$-fold iterate. For a graph
$G$, an \emph{even subgraph} is a set of edges that meets every vertex an even number of times.
With symmetric difference as addition, the even subgraphs form a vector space $\Z(G)$ over
$\mathbb{F}_2$. Its dimension is
\[
  \beta(G)=|E(G)|-|V(G)|+\comp(G),
\]
where $\comp(G)$ is the number of components. We call $\beta(G)$ the \emph{rank} of $G$; other
names are cyclomatic number and cycle rank. A \emph{cycle} is a connected $2$-regular subgraph. We
identify a subgraph without isolated vertices with its edge set. Then a cycle is a nonzero element
of $\Z(G)$, so a graph of rank $r$ has at most $2^{r}-1$ cycles.

Let $\psi(r)$ be the largest number of cycles in a graph of rank $r$.
\begin{itemize}[leftmargin=1.4em]
\item Entringer and Slater~\cite{ES} proved that $2^{r-1}<\psi(r)<2^{r}$ for $r\ge3$, and they
conjectured that $\psi(r)\sim2^{r-1}$.
\item Aldred and Thomassen~\cite{AT} proved that $\psi(r)\le\frac{15}{16}2^{r}$.
\item Alahmadi, Aldred, dela Cruz, Sol\'e and Thomassen~\cite{AACST} proved that a connected
graph with $p$ vertices and more than $2p+3\log(3p)$ edges has at most $2^{r-1}$ cycles.
\item Horak and Kim~\cite{HK} found cubic graphs of rank $r$ with $H(r)$ cycles, where
\begin{equation}\label{eq:H}
  H(2q+1)=2^{2q}+\bigl(2q+\tfrac12\bigr)2^{q}-3q,\qquad
  H(2q+2)=2^{2q+1}+\bigl(2q+\tfrac92\bigr)2^{q}-3q-4 .
\end{equation}
\end{itemize}
For odd rank $r=2q+1$ the graph of Horak and Kim is a ring of $q$ copies of $K_{2,2}$, each joined
to the next by two edges. For even rank, one more edge is inserted across the ring
(Figure~\ref{fig:ring}). We call this graph $R_r$.

\begin{figure}[ht]
\centering
\begin{tikzpicture}
  \def\ri{1.15}\def\ro{1.95}\def\dl{15}
  \foreach \i/\th in {1/18,2/-54,3/-126,4/162}{
    \coordinate (li\i) at (\th+\dl:\ri);
    \coordinate (lo\i) at (\th+\dl:\ro);
    \coordinate (ri\i) at (\th-\dl:\ri);
    \coordinate (ro\i) at (\th-\dl:\ro);
    \draw[cool] (li\i)--(ri\i) (lo\i)--(ro\i) (li\i)--(ro\i) (lo\i)--(ri\i);
  }
  \coordinate (xi) at (90:\ri);
  \coordinate (xo) at (90:\ro);
  \draw[hot] (xi)--(xo);
  \draw[edge] (ri4)--(xi) (ro4)--(xo) (xi)--(li1) (xo)--(lo1);
  \draw[edge] (ri1)--(li2) (ro1)--(lo2) (ri2)--(li3) (ro2)--(lo3) (ri3)--(li4) (ro3)--(lo4);
  \foreach \i in {1,2,3,4}{\foreach \n in {li,lo,ri,ro}{\node[vtx] at (\n\i) {};}}
  \node[vtx] at (xi) {};
  \node[vtx] at (xo) {};
  \node[note, text=warm, anchor=west] at (2.75,1.75) {the further edge (even rank only)};
  \draw[warm, line width=0.4pt] (2.7,1.75) -- (0.12,1.6);
  \node[note, text=accent, anchor=west] at (2.75,-0.55) {a copy of $K_{2,2}$};
  \draw[accent, line width=0.4pt] (2.7,-0.55) -- (1.7,-0.62);
  \node[note, anchor=west] at (2.75,0.6) {two edges join consecutive copies};
  \draw[slate, line width=0.4pt] (2.7,0.6) -- (1.62,0.98);
\end{tikzpicture}
\caption{The graph $R_{10}$: a ring of four copies of $K_{2,2}$ and one further edge. It has rank
$10$ and $H(10)=696$ cycles.}
\label{fig:ring}
\end{figure}

\begin{maintheorem}{the maximum number of cycles}\label{thm:exact}
There is an integer $R$ such that $\psi(r)=H(r)$ for every $r\ge R$.
\end{maintheorem}

In particular $\psi(r)=2^{r-1}+(1+o(1))\,r\,2^{\lfloor(r-1)/2\rfloor}$, so $\psi(r)\sim2^{r-1}$.
The proof gives no value for $R$.

\begin{fineprint}
The theorem is false for small $r$. For example, take the vertices $u_i,v_i$ with
$i\in\mathbb{Z}/26$ and the edges $u_iu_{i+1}$, $u_iv_i$ and $v_iv_{i+10}$. This generalised
Petersen graph has rank $27$, and a direct count by computer gives $67\,959\,973$ cycles, while
$H(27)=67\,325\,913$.
\end{fineprint}

\subsection{The five cases of the proof}

The number of cycles does not change when we delete a vertex of degree at most one, or replace a
vertex of degree two and its two edges by one edge. A vertex of degree at least four can be split
into vertices of degree three, and no cycle is lost. So it suffices to count the cycles of a
connected multigraph $G$ of rank $r$ in which all vertices have degree three
(Lemma~\ref{lem:cubic}). The graphs $R_r$ show that $\psi(r)\ge H(r)$.

For the upper bound there are five cases (Figure~\ref{fig:plan}). A \emph{cyclic cut} is a set of
edges whose deletion leaves two parts that each contain a cycle. The number $H(r)-2^{r-1}$ is of
order $r\,2^{r/2}$, so the five bounds give the theorem (Section~\ref{sec:proof}).

\begin{figure}[ht]
\centering
\begin{tikzpicture}[
  ask/.style={draw=accent, line width=0.7pt, rounded corners=2pt, fill=white,
    font=\sffamily\footnotesize, text width=4.6cm, align=center, inner sep=4.5pt},
  res/.style={fill=accent!7, rounded corners=2pt, font=\sffamily\footnotesize,
    text width=6.3cm, align=left, inner sep=4.5pt},
  ar/.style={-{Stealth[length=4.5pt]}, line width=0.6pt, draw=slate},
  lab/.style={font=\sffamily\scriptsize, text=slate, inner sep=1.5pt},
  node distance=5.2mm]
  \node[ask] (q1) {Does $G$ have a loop?};
  \node[ask, below=of q1] (q2) {Does $G$ have a cyclic cut of at most three edges?};
  \node[ask, below=of q2] (q3) {Does $G$ have a cyclic cut of four edges?};
  \node[ask, below=of q3] (q4) {Does $G$ have a cycle of length less than $D_*$?};
  \node[res, right=13mm of q1] (o1) {at most $2^{r-1}$ cycles\\ \textcolor{slate}{for a trivial reason}};
  \node[res, right=13mm of q2] (o2) {at most $2^{r-1}+O(2^{r/2})$ cycles\\ \textcolor{slate}{Proposition~\ref{prop:threecut}}};
  \node[res, right=13mm of q3] (o3) {a ring, with at most $H(r)$ cycles, or fewer than $2^{r-1}$ cycles\\ \textcolor{slate}{Theorem~\ref{thm:fourcut}, Corollary~\ref{cor:ringmax}}};
  \node[res, right=13mm of q4] (o4) {fewer than $2^{r-1}$ cycles\\ \textcolor{slate}{Theorem~\ref{thm:girth}}};
  \node[res, below=of q4, text width=4.6cm, align=center] (o5) {fewer than $2^{r-1}$ cycles\\ \textcolor{slate}{Theorem~\ref{thm:half}}};
  \foreach \a/\b in {q1/o1,q2/o2,q3/o3,q4/o4} \draw[ar] (\a) -- node[lab, above] {yes} (\b);
  \foreach \a/\b in {q1/q2,q2/q3,q3/q4,q4/o5} \draw[ar] (\a) -- node[lab, right] {no} (\b);
\end{tikzpicture}
\caption{The five cases for a connected cubic multigraph $G$ of large rank $r$. The number $D_*$
is an absolute constant.}
\label{fig:plan}
\end{figure}

\subsection{Outside parts and the forest estimate}

Fix a set $P$ of vertices of $G$. For an even subgraph $\xi$ we call the set of its edges without
an end in $P$ its \emph{outside part} (Figure~\ref{fig:outside}). If a cycle meets $P$, then its
outside part is a proper subset of a cycle, so it is a forest.

\begin{idea}{cycles and outside parts}
Let $W$ be a cycle, $e$ an edge of $W$, and $P=V(W)$. For every cycle $C$ through $e$, the even
subgraphs $C$ and $C+W$ have the same outside part, and it is a forest. By
the \emph{forest estimate} (Theorem~\ref{thm:forest}), the outside part is a forest for at most
$\half+o(1)$ of all even subgraphs when $P$ is small. An edge recursion then gives
$\psi(r)\le(\half+o(1))2^{r}$ (Corollary~\ref{cor:cycles}).
\end{idea}

\begin{figure}[ht]
\centering
\begin{tikzpicture}[x=1cm,y=1cm]
  % W: a hexagon
  \foreach \i in {0,...,5} \coordinate (w\i) at ({60*\i+30}:0.9);
  \draw[cool] (w0)--(w1)--(w2)--(w3)--(w4)--(w5)--cycle;
  % outside vertices
  \coordinate (a) at (2.1,1.05);
  \coordinate (b) at (3.3,0.25);
  \coordinate (c) at (2.1,-1.05);
  \coordinate (d) at (-2.2,0.9);
  \coordinate (e) at (-3.0,-0.1);
  \coordinate (f) at (-2.0,-0.9);
  \draw[edge] (w3)--(d) (w2)--(d) (w3)--(f);
  \draw[edge] (d)--(e)--(f) (a)--(c);
  % the cycle C: w0 - a - b - c - w5 - w0 ; outside part thick
  \draw[hot] (a)--(b)--(c);
  \draw[hot, dash pattern=on 2.2pt off 1.6pt] (w5)--(w0) (w0)--(a) (c)--(w5);
  \foreach \p in {w0,w1,w2,w3,w4,w5,a,b,c,d,e,f} \node[vtx] at (\p) {};
  \node[note, text=accent] at (0,0) {$W$};
  \node[note, text=warm, anchor=west] at (3.45,0.25) {the outside part of the cycle $C$: a path};
\end{tikzpicture}
\caption{A cycle $W$ and a cycle $C$ through an edge of $W$. The edges of $C$ with an end on $W$
are dashed. The other edges of $C$ form its outside part for $P=V(W)$, which is a forest. An even
subgraph whose outside part contains a cycle is not a cycle through $W$.}
\label{fig:outside}
\end{figure}

In the cases 2, 3 and 4 a cycle consists of paths in the parts of $G$. By the forest estimate,
a part with more than half of the possible paths between two of its ends, by a fixed margin, is a
ladder up to a bounded error, and it has few cycles (Section~\ref{sec:poles}). In the case 3
we also close a part by a small tree, a \emph{cap}, and apply the bound $(\half+o(1))2^{r}$ to
the closed part (Lemma~\ref{lem:cap}). In the case 5 the
forest estimate is not sufficient, because we need a bound below $\half$ and its error term is
positive. In this case a second-moment argument gives two disjoint pieces of $G$ with small
boundaries and large girth, and each piece contributes a factor of at most $\frac35$ to the
proportion of forests (Section~\ref{sec:half}).

\subsection{Authorship and verification}\label{sec:status}

KNT is a person. Astra and Fable~5.2 are language models: GPT-6 Astra and Claude Fable~5.2. KNT
directed the work at each step. The forest estimate and its proof are from~\cite{KNT}; KNT
obtained them in September 2026 with Astra, for a problem of Erd\H{o}s on pancyclic graphs (see
also~\cite{Pan}). In November 2026 the two models studied the number of cycles in parallel, in one
shared set of files.

The parts have these origins. Fable~5.2 found the edge recursion, the localisation and layer
argument, the chain of nested cycles, the parallel-link lemma and Lemma~\ref{lem:star}. Astra found
the path structure of Section~\ref{sec:poles}, the gap theorem, the first analysis of cuts, the
count for rings, the argument for bounded girth, and the packing bound and the two girth scales of
Section~\ref{sec:half}. In its first form, Section~\ref{sec:half} gave the bound $\half$ for the
proportion of forests, and the proof of Theorem~\ref{thm:exact} then used an induction on the
rank. Astra improved the bound to a constant below $\half$, and the induction is no longer
necessary. Fable~5.2 shortened the arguments for paths and cuts and wrote this text. After the
formal proof was complete, Astra found the cap of Lemma~\ref{lem:cap}, which removed two cases
from the proof of Theorem~\ref{thm:fourcut}, and the marked set of Case~5 in
Section~\ref{sec:proof}. Fable~5.2 then found that the chain of
Proposition~\ref{prop:chain} also gives the packing bound, and that
Proposition~\ref{prop:loggirth} serves Sections~\ref{sec:forest} and~\ref{sec:half}; this made
Section~\ref{sec:half} shorter.

Theorem~\ref{thm:exact} has a formal proof in Lean~4. The formal statement is
\lean{Erdos1016.ExactMaximum.exactMaximum_caps}: there is a number $R$ such that
\lean{maxCycles r = horakKim r} for every $r\ge R$. Here \lean{maxCycles r} is the largest number
of cycles of a simple graph with cyclomatic number $r$ on a vertex set $\{0,\dots,n-1\}$, a cycle
is counted by its edge set, and \lean{horakKim r} is $H(r)$. A reader who trusts the Lean kernel
needs to read these definitions only. The table lists the formal statements of the main steps.
All of them are proved.

\begin{center}\small
\begin{tabular}{@{}p{6.3cm}p{8.1cm}@{}}
\toprule
\textsf{\textbf{Statement}} & \textsf{\textbf{Declaration}}\\
\midrule
Theorem~\ref{thm:forest} & \lean{Erdos1016.ShortProof.eventually_marked_forest_bound_logSquared}\\
Corollary~\ref{cor:cycles} & \lean{Erdos1016.cycleCountTheorem}; with parallel edges: \lean{EM.numCycles_le_of_large_rank_margin}\\
Lemma~\ref{lem:cubic}, as the upper bound uses it & \lean{EMS.CubicReduction.cubicReductionStatement}\\
Equation~\eqref{eq:glue23} and the count of Lemma~\ref{lem:glue4} & \lean{Net.numCycles_of_three_crossing}, \lean{Net.numCycles_of_four_crossing}\\
Lower bound of Corollary~\ref{cor:ringmax}: for every $r\ge3$ a graph of rank $r$ with at least $H(r)$ cycles & \lean{EM.lowerBoundStatement}\\
Case 1 of Section~\ref{sec:proof} (a loop) & \lean{EM.numCycles_le_horakKim_of_loop}\\
Theorem~\ref{thm:fewcycles} for $\varepsilon=\frac1{1000}$, with Lemma~\ref{lem:reduction} and Proposition~\ref{prop:thin}; Lemma~\ref{lem:gadgets} & \lean{EM.fewCyclesStatement}, \lean{EM.gadget_third_pair}, \lean{EM.gadget_cross_pairs}, \lean{EM.gadget_opposite_pair}\\
Case 2: Proposition~\ref{prop:threecut} with some constant & \lean{EM.threeCutStatement}\\
Theorem~\ref{thm:gap} & \lean{EMS.gapStatement}\\
Lemma~\ref{lem:cap} & \lean{EMS.Shortening.TreeCapCycleBound.exists_thousandth_bound}\\
Lemma~\ref{lem:star}; Lemma~\ref{lem:rigid} and the cases of Theorem~\ref{thm:fourcut}, as arithmetic & \lean{Net.two_mul_numPathPairs_le_of_full}, \lean{EMS.Shortening.FourCutTreeCapArithmetic.common_full_cross_of_tree_caps}\\
In Theorem~\ref{thm:fourcut}: the sides have no bridge; four full pairs in both sides give at most $H(r)$ cycles & \lean{EMS.Shortening.MinimalCutBridgeless.four_cut_shore_bridgeless}, \lean{EMS.ringCaseStatement}\\
Case 3: Theorem~\ref{thm:fourcut} with the rings & \lean{EM.fourCutStatement_caps}\\
Corollary~\ref{cor:nodense}, for the pole alone & \lean{EMS.cyclicShoreStatement}\\
Case 4: Theorem~\ref{thm:girth} & \lean{EM.shortCycleStatement}\\
Proposition~\ref{prop:packing} with the bound $150\,d^{3}(s+1)2^{d}$; Theorem~\ref{thm:half} & \lean{EMS.Shortening.ChainPacking.exists_cutoff}, \lean{EMS.Shortening.ChainForestProbability.exists_cutoff}\\
Case 5 (large girth) & \lean{EM.largeGirthStatement_chain}\\
Upper bound: at most $H(r)$ cycles for large $r$ & \lean{EM.upperBoundStatement_caps}\\
Theorem~\ref{thm:exact}; equation~\eqref{eq:H} for both parities & \lean{EM.exactMaximum_caps}; \lean{EM.maxCycles_formula_caps}\\
\bottomrule
\multicolumn{2}{@{}p{14.4cm}@{}}{\rule{0pt}{2.8ex}Short names: \lean{EM} is \lean{Erdos1016.ExactMaximum}, \lean{EMS} is \lean{Erdos1016.ExactMaximumSupport}, and \lean{Net} is \lean{Erdos1016.BoundaryTrace.Network}.}\\
\end{tabular}
\end{center}

The files are in~\cite{Repo}, on the branch \texttt{feat/cycle-count-bound}. Commit
\texttt{3a80aa9} has a first formal proof, \lean{EM.exactMaximum}. It follows the first form of
this text, which had no cap, marked a shortest cycle in Case~5, and had another proof of
Proposition~\ref{prop:packing}. The table names the second formal proof, which follows this
text; at the date of this draft its files are not in a public commit. The files contain no
\lean{sorry}: Fable~5.2 wrote \lean{Erdos1016/ExactMaximum/}, and Astra wrote
\lean{Erdos1016/ExactMaximumSupport/} and the formal proof of Theorem~\ref{thm:forest}. Astra's
files contain the formal proofs of Section~\ref{sec:full}, of the upper bound for rings, of
Lemma~\ref{lem:cap} with the arithmetic of Theorem~\ref{thm:fourcut}, and of
Section~\ref{sec:half}. The Lean kernel accepts the formal proofs with the axioms
\lean{propext}, \lean{Classical.choice} and \lean{Quot.sound} only.

The kernel check covers the formal statements, and no machine has checked the text of this
paper. A formal proof can differ from the text. In four places the method is different. The
formal proof of Theorem~\ref{thm:forest} bounds the links of two cycles, which is Step~2 of
Proposition~\ref{prop:loggirth}, with Lemma~\ref{lem:smallcuts}. The formal proof of Corollary~\ref{cor:cycles} uses the forest
estimate in its form for a simple graph, which is in~\cite{Pan}. The formal proofs of
Lemma~\ref{lem:ladders} and
Lemma~\ref{lem:thirdport} use two disjoint paths between the ports and exact numbers of paths in
a capped ladder. The formal proof of Lemma~\ref{lem:affine} first deletes the vertices of degree
one, and then a count of the vertices of a multigraph of large girth gives the large layer. In
other places the details differ. For example, the formal number $D_*$ is at least $256$, the
formal cap of Lemma~\ref{lem:cap} has a vertex in the middle of the edge $xy$, the formal proof
of Theorem~\ref{thm:fewcycles} calls a thin layer ordinary if at most one of its two vertices is
on a pair of parallel edges, and the formal proof of Theorem~\ref{thm:girth} allows $12g+2$
ports for the poles $K_p$. So the text can contain errors that the formal proof does not have.
The two models have read the text, and no referee has. Scripts in the repository test the
statements of Sections~\ref{sec:poles} and~\ref{sec:cuts} on small graphs. These tests do not
replace a proof.
